\documentclass[../main.tex]{subfiles}
\begin{document}
%==========================================TIÊU ĐỀ TẬP========================================
\begin{table}[h]
  \begin{adjustwidth}{-1cm}{}
    \begin{tabular}{>{\centering\arraybackslash}p{6cm}|>{\raggedright\arraybackslash}p{12.5cm}}
      \multirow{5}{6cm}
      % Cột bên trái: ÉPISODE và số tập
      {\\[-40pt]
        \flaregothic\fontsize{20pt}{20pt}\selectfont \centering
        \color{eptype}ÉPISODE\color{black}\\
        \burbank\fontsize{72pt}{72pt}\selectfont
        \color{epnum}90
        \\[18pt]
        \flaregothic\fontsize{14pt}{14pt}\selectfont
        \color{epattr}BIENVENUE, \uline{\color{Polmao} Polka} !\\
        % \flaregothic\fontsize{18pt}{18pt}\selectfont
        % \color{epattr}ÉPISODE PERDU
        \color{black}
      }
       &
      % Dòng thứ nhất: tên tập tiếng Nhật
      {\fotskip\fontsize{18pt}{18pt}\selectfont \ruby{猫}{ねこ}になる\ruby{方}{ほう}\ruby{法}{ほう}
      }                                               \\[6pt]
      % Dòng thứ hai: tên tập tiếng Anh
       & {\cafeteria\fontsize{18pt}{18pt}\selectfont How NOT to be a Cat}                                                            \\[4pt]
      % Dòng thứ ba: tên tập tiếng Việt
       & {\vnmsans\fontsize{18pt}{18pt}\selectfont Mèo hay Cáo?}                                                                     \\[8pt]
      % Dòng thứ tư: Nhân vật xuất hiện
       & {\brian{\fontsize{14pt}{14pt}\selectfont Track used: Caliginous Carnival - Ultimate Battle (from {\cafeteria Reflourished})}
         }
      \\[8pt]
      % Dòng cuối cùng: Thời gian diễn ra của tập (thường sẽ là ngày phát hành)
       & {\continuum\fontsize{16pt}{16pt}\selectfont Ngày phát hành: 24/01/2021}                                                     \\[18pt]
    \end{tabular}
  \end{adjustwidth}
\end{table}
%=========================================TRUYỆN CHÍNH========================================
\color{black}

\pov{Hikawa Kuon}

Tôi biết, tôi đã sống sót năm tuần liên tiếp từ các thành viên của \textit{holoForce} với mỗi người có một kiểu phá phách khác nhau làm cho tôi chỉ muốn than trời vì bọn họ quá loạn lạc. Tuy nhiên, tôi cũng đã nói từ trước, đó chỉ là bước tiếp theo, vì sẽ có nhiều thành viên khác sẽ tới văn phòng này.

Và hiển nhiên, ngày đó cũng đã tới. Ba tuần liên tiếp, tôi sẽ phải tiếp đón ba thành viên từ \textit{NePoLaBo}, và họ sẽ là những con người làm cho tôi mệt mỏi hơn đây.

Tuần đầu tiên, chúng ta có gì nào?

Một nghệ sĩ xiếc đến từ một đoàn diễn bí ẩn, luôn tràn đầy năng lượng và không bao giờ để sân khấu nguội lạnh. Cô gái này là một ``chú hề lạc quan'' (mặc dù không tự nhận mình là một chú hề), luôn mang đến tiếng cười nhưng cũng ẩn giấu nhiều tâm tư sâu kín.

Em ấy có một sức hút khó cưỡng: dù thường xuất hiện với vẻ ngoài vui nhộn, ngốc nghếch, hài hước, nhưng thực ra lại là một người có nội tâm rất phong phú. Em ấy có thể giả giọng được bất cứ thành viên nào trong cả công ty \hll\ này và có một quãng giọng khá sâu. Với khả năng làm chủ sân khấu tuyệt vời, em ấy có thể biến bất kỳ khoảnh khắc nào thành một màn trình diễn đáng nhớ.

Không những thế, tính cách của em cũng khá đặc biệt: sôi nổi, náo nhiệt khi ở cạnh bạn bè mà em thân quen nhưng đôi khi lại trở nên trầm lặng khi một mình. Em rất giỏi khuấy động bầu không khí, nhưng cũng có những lúc đắm chìm trong những suy nghĩ riêng. Một nghệ sĩ luôn cười rạng rỡ, nhưng lại hiểu rất rõ cảm giác của những ai đang cô đơn.

Nghệ sĩ xiếc này cũng là một người\dots\ nói nhanh như gió, thậm chí còn ``tổ lái'' sang nhiều vấn đề khác, để rồi quay trở lại với chủ đề chính.

Và dù là một nghệ sĩ xiếc, mục tiêu lớn nhất của em không chỉ là mang đến niềm vui nhất thời. Em ấy muốn một ngày nào đó có thể tổ chức một buổi diễn vĩ đại, nơi mọi người không chỉ cười mà còn thực sự cảm nhận được những điều ý nghĩa từ sâu thẳm trái tim.

Một nghệ sĩ với giấc mơ tỏa sáng.

\dots Hoặc có lẽ, một ``chú hề buồn'' ẩn mình sau lớp mặt nạ rực rỡ.

\chrint

\begin{wrapfigure}[26]{r}{7.5cm}
  \includegraphics[width=7.5cm]{../../../../Image/Book?/polka_model.png}
\end{wrapfigure}

Omaru Polka (\ruby{尾}{お}\ruby{丸}{まる}ポルカ, \ruby{Ô}{Vĩ}-\ruby{ma-ru}{Hoàn}\ Pô(-rư)-ka) có thể coi là trưởng nhóm của Thế hệ thứ Năm. Như đã giới thiệu ở trên, Polka-chan là một nghệ sĩ xiếc gia nhập vào \hll\ với mong muốn trở thành người điều khiển thế giới VTuber.

Một người thích gây ấn tượng với khán giả bằng những màn nhào lộn và cô sống theo phương châm ``Khi bạn đã quyết định làm điều gì đó, hãy theo đuổi đến cùng!''

Polka-chan được biếtđến là một người vừa khờ khạo, vừa táo bạo. Đúng với bản chất của một người điều hành xiếc, em ấy rất năng động và hướng ngoại, và thường nói nhanh hơn cả những người xem bản xứ có thể hiểu được. Cảm xúc của cem ấy cũng thay đổi biến thiên: lúc thì cực kỳ ồn ào và tràn đầy năng lượng, và lúc thì lại im lặng và điềm tĩnh.

Khi lên sóng trực tiếp, Pol-pol cũng rất hỗn loạn và đầy khó đoán: em ấy nói nhanh và lỏng lẻo, nhảy qua nhảy lại giữa nhiều chủ đề khác nhau một cách nhanh chóng. Ngay cả khi có một danh sách các chủ đề để nói, em ấy cũng nhanh chóng chuyển sang các chủ đề liên quan, chỉ quay lại chủ đề ban đầu sau một thời gian dài.

Polka-chan cũng là một người thích đùa nhây do bẩm sinh, vậy nên em ấy cũng là người sẽ bày ra nhiều trò tai quái khiến cho tôi - một người không thích đùa lắm - đau đầu lắm đây. Em ấy là một người thích đùa, nhưng cũng là nạn nhân của những trò trêu chọc và chơi khăm của các thành viên khác, trong đó có tôi.

Pol-pol dễ bị ảnh hưởng thời tiết, nhất là từ\dots\ áp suất khí quyển. Tôi không nói đùa đâu, có lần em ấy đã từng kể với tôi rằng nếu lên sóng vào một ngày có áp suất thấp (thường là trời rét và ẩm). em ấy sẽ không hoạt bát như thường.

Polka-chna là một cô gái cáo fennec (sau đây sẽ gọi là cáo Sa mạc) với đôi tai cáo dài và mềm mại cùng một cái đuôi bồng bềnh. Em có mái tóc vàng dài với những vệt màu hồng và màu đen, những hình các quân bộ bài poker dưới khóe mắt, đôi mắt màu tím với con ngươi hình biểu tượng dường như thay đổi theo cảm xúc mà em ấy thể hiện, thường có hình trái tim hoặc hình ngôi sao.

Pol-pol để một bím tóc tết một bên vai, phần mái được rẽ đôi dài đến giữa mắt và một chiếc nơ màu xanh trên mái tóc hình cánh quạt đặc trưng của mình.

Bộ trang phục em ấy hay diện bao gồm một chiếc mũ hề nghiêng che đi chiếc tai bên trái, kẹp tóc hình chữ X, mỗi kẹp có đính một vật trang trí theo hình các quân bộ bài khác nhau; áo hai dây không tay có khuy nơ theo chủ đề bộ bài và khoét rốn, có cổ áo cánh rời đeo nơ đỏ truyền thống; váy ngắn xếp lớp màu xanh và trắng sọc dọc bồng bềnh có nơ lớn ở sau lưng, tất đùi sọc dọc không đối xứng màu xanh và hồng buộc chéo và giày mũi nhọn màu đỏ không đối xứng; đeo một chiếc găng khuỷu tay đơn buộc chéo nhiều màu đen và đỏ ở tay trái và một dây đeo tay và một vòng cổ tay ở tay kia.

\chrint

Vừa mở màn mà đã có một bản sơ yếu lí lịch đồ sộ sự hỗn loạn này, không biết em ấy sẽ mang đến điều gì cho tôi nữa\dots

\il

Một ngày cận cuối tháng Giêng, tiết trời se lạnh nhưng chưa quá khắc nghiệt. Tôi và Fubu-chan ngồi trên ghế dài trong phòng nghỉ, trò chuyện vu vơ về những loài cáo.

``Cáo có nhiều loại nhỉ, từ cáo tuyết, cáo đỏ cho đến cáo sa mạc,'' tôi nói, chống cằm nhìn ra cửa sổ, rồi liếc về em ấy: ``Mà em có nghĩ mình giống như một con cáo không?''

Fubuki nghiêng đầu, vẫy nhẹ đôi tai cáo trên đầu mình, rồi cười khúc khích. ``Hmm\dots\ nhìn kỹ lại bản thân thì em cũng có hơi hướng giống cáo thật. Nhưng mà một con cáo kiểu gì nhỉ?''

\img{7-8a}

Câu chuyện của chúng tôi chưa kịp đi xa thì bỗng nhiên, một giọng nói tràn đầy nhiệt huyết từ phía sau vang lên: ``Này, này, này! Đây có phải là \hll\ mà em đã từng nghe nhiều qua không?''

Tôi quay đầu lại. Đứng đó là Poika-chan, người đến văn phòng này lần đầu. Cô nàng giơ tay lên trời, cười toe toét như thể đang chuẩn bị bước vào sân khấu biểu diễn.

``Omaru Polka đây xin phép xin một chỗ nha!'' em ấy phấn khích tuyên bố. Nhưng ngay khoảnh khắc ấy, một cơn gió mạnh đột ngột ập đến.

``\textbf{Nhận lấy này!}'' Haato-chan không biết từ đâu xuất hiện, tay bê theo một cái quạt cây cỡ đại, bật số mạnh nhất và chĩa thẳng vào nàng Cáo xiếc. Trong một giây, ``tên lính mới'' bị thổi bay đi như một mảnh giấy, lăn lóc vài vòng trước khi đáp đất trong trạng thái choáng váng.

\begin{figure}[H]
  \centering
  \begin{subfigure}[t]{0.45\textwidth}
    \centering
    \includegraphics[width=8cm]{../../../../Image/Book?/7-8b1.png}
    \caption{Haato bê một chiếc quạt cây, bật số lớn nhất\dots}
  \end{subfigure}
  \quad
  \begin{subfigure}[t]{0.45\textwidth}
    \centering
    \includegraphics[width=8cm]{../../../../Image/Book?/7-8b2.png}
    \caption{\dots làm cho Polka bị thổi đi.}
  \end{subfigure}
\end{figure}

Tôi thở dài, lắc đầu: ``Đừng có bắt nạt người mới chứ em ơi\dots''

Pol-pol loạng choạng đứng dậy, còn tôi đỡ em ấy ngồi xuống ghế. Vẫn chưa hoàn hồn, Polka-chan ngước nhìn tiền bối đầy hoang mang: ``S-Senpai làm như vậy để làm gì chứ?''

Haato-chan cười tủm tỉm. ``Em biết đấy, ấn tượng đầu tiên là điều quan trọng mà!''

Tôi lập tức cau mày: ``Ấn tượng ở chỗ nào chứ? Có mà em bắt nạt ma mới thì có.''

Trong lúc đó, nàng Cáo trắng vẫn đang vui vẻ nhảy điệu Cáo quen thuộc của mình, hoàn toàn không bận tâm đến cuộc hỗn chiến vừa rồi.

\img{7-8c}

Còn Polka-chan, trước khi kịp phản ứng, lại bất ngờ hắt hơi: ``\textbf{Hắt xì!}''

Ngay sau đó, một giọng nói vang lên từ phía sau đầy tinh nghịch: ``A! Em hắt hơi vì em ấy muốn tỏ vẻ lạnh lùng hả?''

Tôi quay lại và thấy Ayame-chan đang đứng khoanh tay, khoái chí nhìn Polka-chan. Còn nàng Cáo xiếc thì rùng mình, run lẩy bẩy như thể lời trêu chọc ấy có thể tạo ra gió lạnh thật sự: ``Ặc! Câu đùa của chị lạnh như đài nguyên vậy à!''

Nàng Quỷ sừng chống hông, cười một cách ranh mãnh, trong khi tôi chỉ biết lắc đầu ngán ngẩm: ``Mấy đứa đừng trêu em ấy nữa.''

\img{7-8d}

Fubu-chan nghiêng đầu nhìn hậu bối của mình. ``Ổn chứ? Nếu thấy mệt thì nằm xuống đi.''

Nói rồi, em ấy nhanh chóng trải sẵn một tấm đệm với chiếc chăn bông dày trên ghế sô-pha.: ``Chúng ta có cả chăn này.''

Tôi nhìn mà giật mình ngạc nhiên. ``Kiếm đâu ra bộ futon nhanh thế!?''

Không ai trả lời tôi. Thay vào đó, Pol-pol ngay lập tức reo lên: ``\textbf{Yahoo!}'' rồi nhảy phóc lên bộ chăn đệm, cuộn tròn như một chú cáo tìm được chỗ trú ấm áp.

Nàng Cáo sa mạc mỉm cười, rúc vào chăn và lẩm bẩm: ``Quả đúng đồng hương với nhau thì đoàn kết với nhau thật á\~{}\~{}\~{}''

Tôi bật cười: ``Ý em là `tuy rằng khác giống nhưng chung một giàn' chứ gì? Không đâu em.''

Sở dĩ tôi nói thế, vì dù là đồng hương, họ vẫn có thể hại nhau. Xã hội hiện đại này, đáng buồn thay, là thế mà. Cáo với nhau chưa chắc đã bênh vực nhau, đôi khi còn đẩy nhau vào tình thế trớ trêu.

Vừa dứt lời, nàng Cáo trắng đột nhiên nghiêng đầu, nở nụ cười tinh quái và buông một câu sét đánh ngang tai: ``Em đang nói gì vậy? Polka là mèo mà, đúng chứ?''

Trong một khoảnh khắc, dường như thời gian ngưng đọng lại. Polka đứng sững, mắt mở to đầy hoang mang. Nhưng chưa dừng lại ở đó, Haato và Ayame cũng nhanh chóng hùa theo:

``Là mèo mà?''
``Là mèo ha!''
``Là mèo\dots''

Cả ba người đồng thanh như thể đang đọc thần chú.

Nàng Cáo xiếc bắt đầu đổ mồ hôi lạnh, ánh mắt lạc đi như rơi vào ảo giác. Em ấy lắp bắp: ``P-Polka\dots\ là mèo ư\dots?''

\img{7-8e}

Tôi bất lực, lắc đầu nhìn cảnh tượng trước mắt: ``Ờm\dots\ không. Cáo sa mạc gần với chó hơn ấy. Mà em ấy dễ bị dắt mũi đến thế sao? Em ấy là chủ xiếc đấy!''

Nhưng đã quá muộn. Polka-chan hoàn toàn chìm vào trạng thái bị thôi miên. Em ấy chớp chớp mắt vài lần, rồi đột nhiên thay đổi tư thế, hai tay giơ lên tạo dáng như một con mèo con.

``Meo, meo\~{} Polka đây là một con mèo nha, meo\~{}''

Cả ba người tiền bối ngay lập tức vây quanh, trêu chọc cô nàng bị đánh lừa dễ dàng đến đáng thương.

Fubuki-chan xoa bụng Polka-chan, miệng khen: ``Ai là cô mèo ngoan nào\~{}?''

Haato cũng xoa bụng theo, còn Ayame thì khoanh tay đứng bên cạnh, thích thú quan sát hai người kia đùa giỡn với ``cô mèo'' mới xuất hiện.

Tôi nhìn cảnh này mà chỉ biết thở dài.

Đúng là nghệ sĩ xiếc có khác. Nhưng thay vì đi pha trò cho người ta thì lại để người ta pha trò cho mình. Rõ ngược đời.

\img{7-8f}

Khi cả ba người vẫn còn đang vui vẻ chơi với ``cô mèo con'' Pol-pol, bỗng nhiên, một giọng nói cất lên đầy uy quyền: ``Oi, oi, oi! Đây là cảnh sát mèo đây!''

Tôi ngẩng lên và thấy cặp đôi Chó-Mèo đang tiến vào, cả hai đều đội chiếc mũ cảnh sát. Tôi nhăn mặt, lẩm bẩm: ``\hl{Sao lại có cả chó ở đây nữa vậy\dots}''

\img{7-8g}

Okayu-chan đặt tay lên hông, nheo mắt nhìn đám người trước mặt, giả vờ nghiêm giọng hỏi: ``Có một con mèo trong tụi này ha? Bọn em đến bắt nó đây!''

Korone-chan lập tức xông lên hùa theo: ``Bắt nó!''

Nàng Cáo xiếc đang trong trạng thái bị thôi miên lập tức hoảng loạn. Em ấy lắp bắp, mặt tái mét: ``C-Con mèo ạ? Trong nhóm chúng em sao?''

Tôi chỉ có thể lắc đầu, hướng mặt về cặp đôi Chó-Mèo kia: ``Không có đâu em ơi. Về đi. Đây chỉ có cáo với quỷ và người thôi.''

Nhưng Polka-chan lại càng run hơn: ``K-Không thể nào. Không có! Ở đây không có ai là mèo hết!''

Tôi chán nản: ``Nói như thế người ta càng nghi ngờ đấy. Mà, em nhớ mình là loài gì rồi mà, đúng không Pol-pol?''

Đúng lúc đó, Fubu-chan bất ngờ bước lên một bước, hai tay chống hông, dõng dạc tuyên bố: ``Không đời nào có đây, meo\dots!''

Và ngay lập tức, Okayu-chan cùng Koro-chan lao đến, hốt thẳng em ấy đi luôn.

\begin{dialogue}
  \speak{Okayu} Tìm thấy nó rồi! Bắt lấy!
  \speak{Korone} Ai-ai-sơ! (Rõ, sếp!)
  \speak{Fubuki} Ể\dots\ Ể?
\end{dialogue}

Polka sững người nhìn Fubuki bị giải đi trước mắt mình, rồi đột nhiên quỵ xuống, tuyệt vọng hét lên: ``\textbf{Fubuki-senpai!!!}''

\img{7-8h}

Em ấy ôm mặt, giọng đầy ân hận: ``Không\dots\ Chị ấy vô tội mà! Tất cả là lỗi của mình\dots''

Tôi khoanh tay nhìn em ấy, nhún vai tỉnh bơ nhận xét: ``Chính ra người mới bị bắt mới là người có lỗi đó\dots''

Không phải than thân trách mình đâu, Polka-chan. Em đang bị roi vào một cú lừa thôi miên từ mọi người đấy.

Haato-chan đập tay lên vai nàng Cáo xiếc, tràn đầy quyết tâm: ``Polka\dots\ Chúng ta cùng giải cứu em ấy nào.''

Ayame-chan cũng gật đầu hăng hái: ``Chúng ta cùng đưa Fubuki-chan trở về nào!''

Tôi lặng lẽ nhìn hai người đó đang giơ nắm đấm lên cao để lấy khí thế. ``\dots Nghe giống lời kêu gọi kháng chiến nhỉ\dots''

Nhưng lời kêu gọi ấy thực ra lại rất hiệu quả. Cả ba cô gái lập tức giơ nắm đấm lên cao và đồng thanh hô vang: ``\textbf{Xung phong!}''

Và thế là chiến dịch giải cứu nàng Cáo trắng Fubuki-chan chính thức bắt đầu.

Tất nhiên, dù không muốn, tôi cũng phải đi theo. Đi một cách miễn cưỡng\dots\ vì tôi là ``chỉ huy trưởng'' của họ mà. Hoặc chính xác hơn, họ cũng sẽ lôi cổ tôi đi.

\ct

Công cuộc giải cứu nàng Cáo trắng này đúng là đầy rẫy những chuyện hài hước đến khó tin.

Ban đầu, chúng tôi quyết định đi hỏi thăm những người xung quanh để tìm tung tích của ``con tin.'' Sora-chan và A-san là hai người đầu tiên mà chúng tôi tìm đến, nhưng cả hai đều lắc đầu, bảo rằng không hề thấy Fubu-chan, hay tung tích của cặp đôi Chó-Mèo kia đâu cả.

Tiếp theo, ``binh nhì'' Polka-chan quyết định dùng đến biện pháp điều tra chuyên nghiệp hơn: em ấy lấy một chiếc kính lúp mà Haato-chan cho mượn rồi cúi sát xuống đất, cố gắng tìm kiếm dấu vết mà ba người có thể đã để lại.

Và rồi\dots\ em ấy đụng ngay một con chó.

Tôi chưa kịp phản ứng thì đã thấy nàng Cáo sa mạc chăm chú soi xét sinh vật bốn chân trước mặt, làm như nó là một manh mối quan trọng. Nhưng có vẻ con chó kia không thích bị nhìn chằm chằm như thế, và thế là nó gầm gừ một tiếng, rồi đột nhiên lao đến đuổi theo em ấy.

Pol-pol hét lên hoảng loạn, chạy vòng vòng như con thỏ bị săn đuổi. ``\textbf{T-Tôi không phải mèo đâu mà rượt tôi!?}''

Tôi đưa tay xoa trán: ``Thế em không định giải thích về bộ dạng em hiện giờ à?''

Nhưng chưa hết, sau khi thoát khỏi con chó (bằng cách trèo lên cột điện), nàng Cáo sa mạc lại vô tình vấp phải một miệng cống mở\dots\ và thế là em ấy rơi thẳng xuống đó.

Tiếng hét thất thanh của Polka vang vọng khắp con phố.

Tôi thở dài, trong khi hai ``binh nhất'' Haato-chan và Ayame-chan thì hoảng hốt chạy đến miệng cống, gọi lớn: ``Polka-chan!? Em ổn chứ!?''

Tiếng rên rỉ từ dưới vọng lên: ``\dots Em có thể ngửi thấy quá khứ\dots''

\begin{dialogue}
  \speak{Haato} \dots Ờm, em ấy còn nói được nghĩa là còn sống.
  \speak{Ayame} Nhưng chắc chắn là có mùi lắm đấy, yo.
  \speak{Kuon} \dots Ừ. Mau tìm cách kéo em ấy lên nào.
\end{dialogue}

Và đúng như dự đoán, khi chúng tôi kéo Polka lên khỏi cống, một mùi kinh khủng bốc ra, khiến ai nấy đều phải bịt mũi.

Thế là không còn cách nào khác, chúng tôi phải dẫn em ấy đến một trạm rửa xe gần đó để gột sạch thứ mùi khó chịu kia.

\begin{figure}[H]
  \centering
  \begin{subfigure}[t]{0.45\textwidth}
    \centering
    \includegraphics[width=8cm]{../../../../Image/Book?/7-8i1.png}
    \caption{Biệt đội đã hỏi A-chan và Sora-chan, nhưng chỉ nhận về cái lắc đầu không biết từ họ.}
  \end{subfigure}
  \quad
  \begin{subfigure}[t]{0.45\textwidth}
    \centering
    \includegraphics[width=8cm]{../../../../Image/Book?/7-8i2.png}
    \caption{Polka sử dụng kính lúp để tìm tung tích của ba người thì chạm trán với con chó.}
  \end{subfigure}
  \quad
  \begin{subfigure}[t]{0.45\textwidth}
    \centering
    \includegraphics[width=8cm]{../../../../Image/Book?/7-8i3.png}
    \caption{Con chó rượt đuổi Polka đi, theo sau đó là cô nàng bị rơi xuống cống.}
  \end{subfigure}
  \caption{Tóm tắt ``hành trình'' giải cứu Fubuki của ba cô gái và Kuon đầy gian truân nhưng hài hước.}
\end{figure}

\ct

Cuối cùng, sau khi ``vượt qua bao khó khăn hoạn nạn'' - ít nhất là theo lời của ba cô gái - chúng tôi cũng tìm được hai kẻ bắt bớ kia.

Okayu-chan và Korone-chan đang đứng ở một phía sân thượng đeo cặp kính râm như những người thuộc xã hội đen, trong khi chúng tôi đứng ở sân thượng đối diện. Ở giữa là một thanh sắt nhỏ, nối hai bên lại với nhau.

\img{plot}

Nàng Mèo khoanh tay, nhìn chúng tôi với ánh mắt đầy thách thức: ``Nếu ai đó trong số các người có thể tới đây được\dots''

Nàng Khuyển tiếp lời ngay sau đó: ``\dots bọn em sẽ trả lại Fubuki-senpai lại.''

Cả nhóm chúng tôi im lặng nhìn nhau. Rõ ràng là chúng tôi đang bị buộc phải tham gia một thử thách không hề đơn giản.

Tôi nhìn sang thanh gỗ hẹp nối hai tòa nhà, lòng không khỏi nhíu mày. Đây là con đường duy nhất để băng qua, và chỉ cần sơ suất một chút thôi là có thể ngã xuống dưới. Phải là người cực kỳ khéo léo mới có thể làm được.

Mà nhắc đến khéo léo thì\dots\ chắc chỉ có người đó thôi nhỉ\dots

``Cứ để đó cho em! Em làm được!'' cô nàng Cáo xiếc reo lên đầy tự tin.

Haato-chan phấn khích vỗ tay: ``Thế mới là cô gái xiếc của chúng ta chứ!''

Ayame-chan cũng phụ họa theo, giơ tay làm động tác cổ vũ: ``Thuyền trưởng xiếc của chúng ta kìa\~{}! Fooo\~{}!''

Nghe hai cô gái kia khen, Polka-chan lập tức phổng mũi, khoanh tay lại đầy tự hào. Tôi thấy thế liền xen ngang, kéo em ấy trở về thực tại:  ``Khen gì thì khen sau, bây giờ tập trung đưa Fubuki-chan về trước đã!''

Nhưng em ấy không những không tụt hứng mà còn dõng dạc đáp: ``Không phải đâu, vì Polka đây là mèo mà!'' Nói rồi, em ấy khéo léo giữ thăng bằng trên thanh sắt hẹp. ``Mấy cái này với em thì nhằm nhò gì!''

\dots Rồi cuối cùng em ấy có phải là mèo hay không đây?

\img{7-8j}

Ở bên kia, Fubu-chan kêu lên lo lắng: ``\textbf{Polka à! Đừng lại đây, nguy hiểm lắm!}''

Tôi lắc đầu, lẩm bẩm: ``\hl{Giờ mình còn chẳng biết em ấy khéo vì nghĩ mình là mèo hay vì chính em ấy là nghệ sĩ xiếc nữa\dots}''

Nhưng chưa kịp nghĩ nhiều, Polka đã chạy vun vút trên thanh sắt hẹp, tiến đến cặp Oka-Koro, với những lời cổ vũ cuồng nhiệt từ đồng đội phía sau. Cô nàng Cáo xiếc nhanh chóng lấy ra một quả pháo - không biết từ đâu ra - rồi ném thẳng về phía đối thủ. Một tiếng nổ vang trời khiến Okayu và Korone giật bắn mình, ngã sõng soài ra đất.

\img{7-8k}

Lợi dụng cơ hội, tôi lập tức lao tới, hét lớn: ``\textbf{Đứng lại, bọn kia!}''

Okayu-chan và Koro-chan hốt hoảng định bỏ chạy, nhưng tôi đã tính toán trước đường đi của họ. Chỉ sau vài bước di chuyển, tôi chặn đầu được cả hai và vô tình đánh ngất luôn trong lúc bắt giữ.

\dots Cơ mà, tại sao tôi lại đánh ngất hai em ấy? Có vẻ như động tác này hơi thừa\dots\ E hèm, thôi bỏ đi.

Nàng Quỷ sừng và nàng Song tính từ phía bên kia phấn khích reo lên: ``Em ấy làm được rồi kìa!''

Lúc này, Polka-chan vươn tay ra với ``con tin'' Cáo trắng, nở một nụ cười chiến thắng: ``Ta quay về nào, Fubuki-senpai.''

Mọi chuyện tưởng chừng đã ổn thỏa và đã có thể kết thúc một cách êm đẹp, nhưng\dots\

Haato-chan bỗng nhiên gọi lớn, giọng có chút lo lắng: ``Hai người lại đây nhanh lên nào!''

Ayame cũng giơ tay làm hiệu, ánh mắt có vẻ cảnh giác.

Tôi chợt có một linh cảm xấu\dots

\pov{-}

Polka cõng Fubuki, cẩn thận bước trên thanh sắt hẹp nối hai tòa nhà. Cô nàng Cáo sa mạc có vẻ rất tự tin vào khả năng giữ thăng bằng của mình, nhưng ánh mắt Fubuki lại hiện lên chút băn khoăn.

\img{7-8l}

Cô nàng Cáo trắng ngập ngừng, rồi chậm rãi lên tiếng: ``Polka à, chị\dots''

Polka hơi nghiêng đầu: ``Sao thế ạ?''

Fubuki mím môi, rồi cuối cùng cũng quyết định nói ra: ``Chị phải xin lỗi em vì đã lừa dối em. Em không phải là một con mèo, em là một con cáo sa mạc.''

Câu nói ấy như một gáo nước lạnh dội thẳng vào đầu Polka. Cô khựng lại, mắt mở to đầy sửng sốt.

Em\dots\ không phải mèo sao?

Khoảnh khắc chấn động ấy khiến cô mất tập trung, và rồi\dots\

Trượt chân!

Cả hai người rơi xuống.

Ayame và Haato hét lên hoảng hốt: ``\textbf{Polka! Fubuki-senpai!}''

Nhưng trong khi rơi xuống, hai người họ vẫn mỉm cười với nhau. Không có chút sợ hãi hay thảo mai nào trong ánh mắt họ nữa.

Họ\dots\ đã thực sự trở thành bạn của nhau. Một điều rất thiêng liêng và hiếm thấy khó tìm ở thời đại này.

%==========================================TRUYỆN PHỤ=========================================
\omk

\pov{Omaru Polka}

Ôi chà\dots\ Vui thật đấy.

Fubuki-senpai cuối cùng cũng đã làm hòa với mình.

Mình cứ tưởng là sẽ rơi thẳng xuống đất cơ. Nhưng ngay khoảnh khắc sắp chạm đất, một thứ gì đó mềm mềm, êm êm đỡ lấy mình.

\img{7-8m}

*BỘP!*

Mình tiếp đất một cách không nhẹ nhàng lắm. Dù sao thì rơi từ trên cao xuống cũng chẳng dễ chịu gì. Cả người mình ê ẩm, nhưng may là không bị thương gì nặng.

Fubuki-senpai ở ngay bên cạnh, cũng đang nằm bẹp trên cái phao đệm cứu hộ siêu to. Chị ấy nhìn mình và cười nhẹ. Một nụ cười không còn trêu chọc, mà là một nụ cười thực sự\dots\ như một người bạn.

Lúc đó, mình mới chợt nhận ra: mình không còn cảm thấy bực bội hay ganh tị với chị ấy nữa.

Bọn mình\dots\ đã thực sự trở thành bạn rồi nhỉ?

\dots

Haato-senpai và Ayame-senpai chạy tới, thở hổn hển trong sự nhẹ nhõm:

\begin{dialogue}
  \speak{Haato} May quá trời\dots
  \speak{Ayame} Polka-chan giỏi lắm á!
\end{dialogue}

Mình gượng dậy, chưa kịp nói gì thì bỗng có một giọng nói phía sau cất lên, nghe vừa bực bội vừa mệt mỏi: ``Còn anh thì sao hả!?''

Mình quay đầu lại, thấy anh Tổng ló mặt ra từ phía sau tấm phao, trông có vẻ kiệt sức.

Bốn đứa bọn mình cùng đồng thanh: ``\textbf{Kuon-san/senpai!}''

Anh ấy thở hắt ra, chống tay vào hông, chán nản lắc đầu: ``Anh biết ngay kiểu gì cũng sẽ có chuyện xảy ra mà, thế nên anh đã tính trước rồi.''

Mình tròn mắt ngạc nhiên: ``Sao anh biết được hay vậy?''

``Anh là quản lý của mấy đứa nên biết chứ. Không lẽ mấy đứa muốn anh với cả công ty này trả viện phí à!?'' anh ấy gắt lên, đáp lại. ``Không có đâu em ơi.''

Mình bật cười: ``Em không nghĩ anh cũng có tài quan sát ghê ha!''

``Cái này bình thường mà. Có gì phải ngạc nhiên đâu?'' Kuon-senpai thở dài, bình tĩnh lại đôi chút.

Ayame-senpai chen vào, chỉ tay về phía hai kẻ vẫn còn nằm ngất xỉu sau cú đánh vô tình của anh Tổng: ``Cả cái cách mà anh bắt sống được Okayu-chan với Korone-chan nữa!''

Anh ấy nhún vai, lắc đầu: ``Cái này thì vô tình anh biết thôi mà - Mà khoan, tại sao mấy đứa đây lại chơi trò cảnh sát với bọn này nhỉ?''

Haato-chan khoanh tay lại, bĩu môi: ``Bọn nó thích đùa giỡn thôi anh.''

Kuon-senpai chợt chống tay lên cằm, suy nghĩ một hồi laua.``Chơi trò cảnh sát à\dots''

\dots

\dots

Ánh mắt lóe lên một ý tưởng gì đó. Anh nhếch mép cười: ``Hai đứa ấy muốn chơi bọn anh chứ gì? Vậy để bọn này chơi lại.''

Bốn đứa bọn mình nhìn nhau, rồi quay lại nhìn cậu ấy với vẻ ngạc nhiên. Anh ấy muốn dùng cách mà hai chị ấy đuổi bắt khắp thành phố như vừa sao?

Hay, anh ấy lại có bài gì nữa?

Kuon-senpai chỉ về phía mình và Fubuki-senpai, nói với giọng đầy ẩn ý: ``Anh thấy hai người đã trở thành `đồng chí' của nhau rồi đó.''

Cáo trắng-senpai cười khì, còn mình thì ngượng đỏ mặt: ``Anh cứ nói thế nào ấy chứ\dots''

Cả ba đứa còn lại cùng quay sang nhìn cặp chó-mèo đang bất tỉnh, vừa tò mò với bài vở anh ấy sẽ tung ra, vừa hơi ái ngại cho hai chị ấy.

Giọng của anh ấy bỗng trầm xuống, mang theo một chút gì đó khiến sống lưng mình lạnh toát: ``Nợ máu phải trả bằng máu.''

Mình nuốt nước bọt.

``Ta sẽ dùng chiêu đập ông gậy lưng ông!''

Chúng mình bắt đầu hiểu ra ý đồ của anh ấy. Nhưng rồi, senpai cười một cách đầy nguy hiểm: ``Nhưng mà anh muốn tất cả bốn đứa sẽ thực hiện trò chơi này.''

Bốn đứa bọn mình nhìn nhau đầy háo hức. Ai mà ngờ được cậu con trai kia lại có một cái đầu cũng tinh quái ra phết nhỉ? Nếu Okayu-senpai và Korone-senpai thích đùa giỡn như vậy, thì bọn mình cũng sẽ đáp lại theo đúng tinh thần đó!

Kuon-senpai chống tay vào hông, khẽ bẻ khớp ngón tay, ánh mắt sắc bén hơn thường lệ: ``Để xem trình độ xiếc của Polka-chan như thế nào!''

Hả? Sao tự dưng mình lại bị lôi vào đây thế!?

``K-Khoan đã! Em tưởng là mình chỉ đóng vai `chiến binh' thôi chứ!?''

Anh ấy nhếch mép: ``Không đâu, em sẽ là một trong những `trùm cuối' đấy.''

Mình chỉ biết há hốc mồm trong khi hai người đồng đội mình kéo chị Khuyển đi một hướng, còn Fubuki-senpai và anh Tổng thì lôi chị Mèo về hướng khác. Hai nàng chó-mèo đã được tách riêng hoàn toàn.

Mình cười gượng. Không biết đây là thử thách hay là trừng phạt đây\dots

Và trò chơi chính thức bắt đầu.

\ct

Khi Korone tỉnh dậy\dots

``Hửm?''

Cô nàng Cún ngơ ngác mở mắt, nhìn xung quanh, nhưng lại đụng trúng vào một thùng các-tông chất đầy xung quanh. Một căn phòng tối om, chỉ có vài tia sáng le lói hắt qua khe cửa, đủ để chị ấy nhận ra mình đang ở đâu. Mùi giấy, mùi mực in và cả\dots\ mùi thức ăn thừa!?

``\dots Đây là đâu!?''

Korone bật dậy, lảo đảo nhìn quanh. Từng chồng tài liệu chất đống khắp nơi, một cái bàn làm việc lộn xộn, và một vài hộp đồ ăn nhanh vứt ngổn ngang. Không thể nhầm được: đây là nhà kho văn phòng!

Korone-senpai nhanh chóng kiểm tra túi áo. Không có điện thoại. Cũng chẳng có bộ đàm.

Chị ấy thử kéo cửa ra\dots\ nhưng nó đã bị khóa. ``Ể!? C-Cái gì vậy trời''

``Tại sao mình lại bị nhốt trong này chứ!?''

Đúng lúc đó, một giọng nói vang lên từ chiếc loa nhỏ trên tường. Một giọng nói vừa quen vừa lạ. ``Chào mừng đến với thử thách sinh tồn, Inugami Korone.''

Korone-senpai giật mình: ``Ơ!? Ai đó!?''

Giọng nói đó tiếp tục, lần này kèm theo một chút chế giễu: ``Muốn cứu Okayu hả? Trước tiên, phải thoát khỏi đây đã.''

Chị Khuyển nhíu mày, nắm chặt tay. ``Okayu đang gặp nguy hiểm ư!? Ngươi đang để cậu ta ở đâu!?''

``Cái đó thì ngươi phải tìm thôi. Một chỗ khá là xa lạ đấy.'' giọng nói ấy - thực chất là Kuon-senpai - đáp lại.

Chị ấy nhìn quanh căn phòng một lần nữa, lần này ánh mắt nghiêm túc hơn. Nếu đây là một trò chơi, thì cô nhất định sẽ thắng.

Trò chơi `giải cứu Okayu' đã chính thức bắt đầu!

\il

Cùng lúc đó, ở một nơi hoàn toàn khác\dots

Một tiếng động lớn vang lên. Cánh cửa sắt nặng nề mở ra, để lộ một nhà xưởng rộng lớn.

Okayu-senpai chớp mắt vài lần, cố gắng nhìn quanh.
``\dots Cái gì vậy trời?''

Mùi dầu máy và sắt thép xộc thẳng vào mũi. Những cái giá đỡ cao chất đầy linh kiện kim loại. Xa xa, những tấm bạt che phủ các cỗ máy lớn, khiến nơi này trông như một bối cảnh phim hành động.

Chị ấy bị trói vào một chiếc ghế gỗ, xung quanh toàn là dụng cụ cơ khí, mô hình máy móc và những mảnh linh kiện lỉnh kỉnh.

Chị thở dài, lười biếng tựa đầu vào lưng ghế, lẩm bẩm: ``Tự nhiên bị bắt cóc qua tận đây\dots\ Mà, đây là đâu đây?''

Trên cao, một chiếc loa cũ kỹ cất giọng: ``Xin chào, Nekomata Okayu. Chào mừng đến với `pháo đài' của bọn ta.''

Okayu-senpai đảo mắt xung quanh, khóe môi cong nhẹ. ``À há\dots\ Vậy là mấy người định chơi trò `con tin' với tôi à?''

Một tràng cười khúc khích vang lên từ phía bên kia nhà xưởng. ``Chính xác! Nếu muốn thoát khỏi đây, ngươi sẽ phải đối đầu với bọn ta!''

Cánh cửa xưởng bỗng mở ra. Tôi chầm chậm bước vào, hé lộ dần trước mặt chị Mèo. Đôi mắt màu tím hình trái tim của tôi ánh lên vẻ thích thú. ``Polka ở đâu? Polka ở đây này!''

Okayu híp mắt. ``\dots Polka, đúng không?''

Mình ngả người ra sau, cười toe toét: ``Đúng rồi đó, Okayu-senpai! Chào mừng đến với vở diễn đặc biệt của bọn em!''

Ở phía xa, Fubuki-senpai, Haato-senpai và Ayame-senpai đã vào vị trí. Còn bên trong phòng điều hành (vốn dĩ là phòng nghỉ, nhưng được anh ấy canh tân lại), Kuon-senpai - người soạn kịch bản điên rồ này - đứng khoanh tay, theo dõi toàn bộ tình hình.

Mọi thứ đã sẵn sàng.

Trò chơi săn đuổi đã chính thức khởi động.

\il

Và đó chính là cách mà Omaru Polka mình đây - thành viên đầu tiên của thế hệ thứ Năm \hll\ - ra mắt.

Nhưng câu chuyện của \textit{NePoLaBo} chúng mình chưa kết thúc ở đây.

Tuần sau, một cô gái khác sẽ bước vào sân khấu.

Một người có mái tóc xanh như băng giá. Một ánh mắt lấp lánh như tinh thể pha lê. Và một nụ cười dịu dàng nhưng ẩn chứa sức mạnh không ngờ.

Hãy chờ xem. Câu chuyện của cậu ấy\dots\ sẽ bắt đầu sớm thôi. Kuon-senpai, liệu anh có thể đối mặt với cô nàng Yêu tinh tuyết băng giá này được hay không?

Cậu ấy sẽ chờ anh đấy. Chúc anh may mắn!

\end{document}